\cotefile{15}
\scanpage{184}\mspage{73}
a homomorphism of restriction of a function on $\mathscr{E}_{\unsure{\varphi|\omega}}$. We obtain from (16):
\handwritten{\[
0\to O_V(-\mathscr{E}_{\varphi})[nE]\to O_V[nE]\xrightarrow{r'}O_{\mathscr{E}_{\varphi}}[nE]_{\mathscr{E}_{\varphi}}\to 0
\tag*{(17)}
\]}
here $r'$ is such: if $P\notin\mathscr{E}_{\varphi}$, then it is a zero homomorphism; if $P\in\mathscr{E}_{\varphi}$ and $\{\alpha^{(i)}\}\in O_V[nE](P)$, then $r'\{\alpha^{(i)}\}=\{r\alpha^{(i)}\}$. To the sequence (17) corresponds such exact sequence:
\handwritten{\[
\begin{aligned}
&H^0(V,O_V[nE])\xrightarrow{r^{\unsure{*}}} H^0(\mathscr{E}_{\varphi},O_{\mathscr{E}_{\varphi}}[nE]_{\mathscr{E}_{\varphi}})\to H^1(V,O_V(-\mathscr{E}_{\varphi})[nE])\to\\
&\to H^1(V,O_V[nE])\to H^1(\mathscr{E}_{\varphi},O_{\mathscr{E}_{\varphi}}[nE]_{\mathscr{E}_{\varphi}})
\end{aligned}
\tag*{(18)}
\]}

Let $n(\mathscr{E}_{\varphi})$ be such that, at $n\geq n(\mathscr{E}_{\varphi})$ for $\mathscr{E}_{\varphi}$ and $[nE]_{\mathscr{E}_{\varphi}}$ satisfy the requirements $1')$ and $2')$ (See p. 202). The existence of such $n(\mathscr{E}_{\varphi})$ follows from the condition of non-degeneracy of $[E]_{\mathscr{E}_{\varphi}}$ over $\mathscr{E}_{\varphi}$. Put
\handwritten{\[
n_{\mathscr{E}}=\max(n'',\,n(\mathscr{E}_{\varphi})).
\]}

Then, at $n\geq n_{\mathscr{E}}$, from (15), (18) and from the fact that, at $n\geq n(\mathscr{E}_{\varphi})$
% archive sheet number: 183
\scanpage{185}
De plus, on a une structure \unsure{polarisée|polaris.|privilégiée}\inserted{ $\pi$} sur $G_{\mathbb{R}}^{\circ}$, qu'on peut \deleted{\unsure{défi|définir}} supposer définie par un hom \deleted{$\unsure{\mu_2\to G}$} $i:\mu_{2\,\mathbb{C}}\longrightarrow G_{\mathbb{C}}^{\circ}$ (comme on voit en envoyant un tore \inserted{\unsure{provenant d'un $s$|prenant un sens}} dedans). D'autre part, $G(s')_{\mathbb{C}}$ est muni de $i(s'):G_{m\,\mathbb{C}}\to G(s')_{\mathbb{C}}$, et \unsure{on|qu'on} obtient par conséquent
\[
\unsure{v_{\,\xi}|v_{1}|\psi}(s')=u(s')_{\mathbb{C}}\circ\bigl(i(s')\mid\mu_2\bigr):\mu_{2\,\mathbb{C}}\longrightarrow G_{\mathbb{C}}^{\xi(s')}
\]
\deleted{D'autre part,} Ceci posé, il faut que $\pi^{\xi(s')}$ (polarisation de $G_{\mathbb{R}}^{\xi(s')}$ déduite de $\pi$ sur $G$ par transport) soit celle définie par les $\mu_{2\,\mathbb{C}}$. Par conséquent, pour $s'$ \unsure{variable} de façon que $\xi(s')$ \unsure{reste|soit} constant, les polarisations \inserted{de $G_{\mathbb{R}}^{\xi(s')}$} associées aux différents $v(s')$ doivent être identiques -- ce qui donne une condition non tautologique sur la distribution des \deleted{\unsure{$v$}} $u(s')$. [\unsure{Noter|Notons} cependant que l'ens des \deleted{\unsure{classes de}} polarisations d'un \unsure{grpe|gpe} réel $H$, défini en termes d'hom $\mu_{2\,\mathbb{C}}\to H_{\mathbb{C}}$, \margin{? à vérifier} est \uline{fini}, [si $G$ est commutatif en tous cas !], donc la restriction n'est pas tellement draconienne.

Enfin, conformément aux conjectures de Tate magnifiées, la connaissance des $u(s')$ \unsure{permet} de retrouver la structure d'effectivité dans la catégorie
% archive sheet number: 184
\scanpage{186}\mspage{72}
Since
\handwritten{\[
O_V(-E_{\varphi})[nE]\cong O_V[(n-1)E]
\]}
then, having taken $n''\geq n'$, we obtain from (14) that, at $n\geq n''$
\handwritten{\[
\begin{aligned}
&\dim H^1(V,O_V[nE])=\\
&=\dim H^1(V,O_V(-\mathscr{E}_{\varphi})[nE])=\dim H^1(V,O_V[(n-1)E]).
\end{aligned}
\tag*{(15)}
\]}

\margin{p. 208}It is evident that $\unsure{\varphi|\omega}^{n}=\{\varphi_i^{n}\}$ is a cross-section of the group $H^0(V,O_V[nE])$. Therefore, the common zeros of the elements of the group $H^0(V,O_V[nE])$ are contained among the zeros (of) $\varphi^n$, that is, all of them lie on the subvariety $\mathscr{E}_{\varphi}$.

Let us show that, at $n\geq n_{\mathscr{E}}$ there is no common zeros of the cross-sections of $H^0(V,O_V[nE])$ on $\mathscr{E}_{\varphi}$. Consider an exact sequence
\handwritten{\[
0\to O_V(-\mathscr{E}_{\varphi})\to O_V\xrightarrow{r}O_{\mathscr{E}_{\varphi}}\to 0
\tag*{(16)}
\]}

The homomorphism $r$ is described in the following way: if $P\notin\mathscr{E}_{\varphi}$, then $r$ in the stalk over this point is a zero homomorphism; if $P\in\mathscr{E}_{\varphi}$, then $r$ over $P$ is
% archive sheet number: 185
\scanpage{187}
des $G_{\mathbb{R}}$-modules de dim finie sur $\mathbb{R}$. Le cône qui exprime cette effectivité est \uline{engendré} par les \deleted{$\unsure{u(s')|v(s')}$} rayons définis par les $v(s')$. \unsure{C'est} un fait qui exprime \inserted{sous forme ``analytique''} le fait que les \unsure{frobénius|frobenius|Frobenius} $f(s')\in G_{\mathbb{R}}^{\xi(s')}(\mathbb{R})$ ont la vertu que leurs valeurs propres, dans \unsure{tte} représentation de $G_{\mathbb{R}}^{\xi(s')}\deleted{(\mathbb{R})}$ -- et en fait \unsure{$=$|même} de $G_{\deleted{\mathbb{R}}\mathbb{C}}^{\xi(s')}$ -- sont des \uline{entiers} algébriques, ce qui est une condition arithmétique assez draconienne !
% archive sheet number: 186
\scanpage{189}
\uline{Théorie de Galois des Motifs}

\uline{\unsure{Généralités|Généralites}}
% archive sheet number: 188
\scanpage{190}
Foncteurs fibres motiviques et \unsure{grpe|gpe|gp} fondamental motivique
% archive sheet number: 189
\scanpage{191}
$X$ préschéma connexe, \deleted{\unsure{(réduite, ce qui ne restreint p|(réduction, ce qui ne restreint p|évidemment, ce qui ne restreint p|réduit, ce qui ne restreint pas}}

$1^\circ)$ $X_{\mathrm{réd}}$ de car $p>0$. Alors dans le \unsure{pro-torseur|protorseur} fondamental $(T_\alpha)$, il n'y a pas en général de sections pour les $T_\alpha$. \deleted{\unsure{On}} On dispose uniquement, \deleted{de façon canonique}

(i) des $\xi_{\ell}^{(\alpha)}\in T_\alpha(\mathbb{Z}_{\ell})$, $l\ne p$, provenant sans doute de $\xi_{\ell}\in T_\alpha(\widetilde{\mathbb{Z}}_{\ell\mathbb{Z}})$
% slanted
\margin{Coh. $l$-adique}
\margin{hensélisé}

\diagram{\begin{tikzcd}[ampersand replacement=\&,column sep=0pt,row sep=3.3em]
\text{(ii) d'un }\xi_p^{(\alpha)}\in T_\alpha( \& \unsure{\mathbb{W}|W}_p(k)\arrow[dr,swap,"{\text{``DeRham-Hodge''}}"] \& )\quad\text{si } \deleted{X\ \unsure{\text{de t.f. sur}|\text{déf.}}}\inserted{\text{$k$ est un}}\ \inserted{\text{corps parfait sur $X$}}\,\deleted{\text{le corps parfait }k} \\
\text{(iii) }\bigl[\,?\quad\text{\uline{réduit}}\ \xi_p^{(\alpha)\,\unsure{0|\circ}}\in T_\alpha( \& \deleted{\unsure{\mathbb{W}}}\ \& k\,)\ \bigr]
\end{tikzcd}}
% slanted
\margin{Coh. $p$-adique. Structure supplémentaire : opérations \fbox{\uuline{$F$} et \uuline{$V$}} $1/2$ \uuline{linéaires} \uuline{\unsure{reliées|reliant|relativement}} : \unsure{Frobenius uni|Frobenius univ}}
% slanted
\margin{Coh. Hodge-de Rham Mais ici, il faut \deleted{sans doute} \unsure{donc|faut} aller dans la catégorie $\overline{\unsure{M|\mathcal{M}}}$ dérivée. Tenir compte aussi de la structure \uline{filtrée} en Hodge-De Rham}

[$\xi_p$ provenant sans doute d'un $\xi_p\in T_p(\overset{\sim\!\sim}{\mathbb{Z}}_{p\mathbb{Z}})$]\margin{hensélisé strict.}

Si $X=\operatorname{Spec}\mathbb{F}_q$, il y a en général des \deleted{\illegible{1}} obstructions non triviales dans des $H^2(\mathbb{Q},G)$ \ldots

\deleted{\unsure{On}} On utilise des ``foncteurs fibres'' qui induisent sur les motifs effectifs de poids $0$ le foncteur canonique. Alors les \unsure{objets|élts|d pts|points} de $T_\alpha(\mathbb{Z}_\ell)$ sont sans doute uniques \inserted{à isom près}, (th. de Lang), de \unsure{$=$|même} ceux de $T_\alpha(\deleted{\unsure{\mathbb{Z}}}\mathbb{W}_p(k))$ si $k$ fini (id), \inserted{plus gén.} pour un $k$ parfait de dim \unsure{ch|coh|coh.} $1$ (?). En général, il faut s'attendre que $T_\alpha(k)$ \deleted{\unsure{soit un}} ait plus d'un élément à iso près, mais il est possible que $T_{\deleted{\unsure{\alpha|\infty}}}(k)=\varprojlim T_\alpha(k)$ n'a qu'un élément à iso près, celui donné par la théorie de DeRham-Hodge.
% slanted
\margin{\unsure{relation} avec les \unsure{problèmes} galoisiens ?? Il doit être \uline{\unsure{équivalent|équivariant}}, cf \uline{Manin}}
% archive sheet number: 190
\scanpage{192}
ALGÈBRES de LIE SEMI-SIMPLES \mspage{(9)}\hfill Tribu n\textsuperscript{o} 59

Pas de noms propres à la p. 55 ; Chevalley et Harish on parlera longuement de dans la note historique. Dans la prop. 7 appeler les gens donnés $X'$ et $Y'$. Au lieu du ``tout de suite'', (\unsure{l.2|1.2}, p. 56), dire pourquoi, (ça se propage à partir des racines simples).

Utiliser le lemme sur les applications $R\longrightarrow\text{groupe}$ au milieu de la p. 56. Ensuite poser $X_\alpha=u_\alpha^{-1}X'_\alpha$, (alors $t_\alpha=1$). Pas de nom à la p. 57.

\uline{Exer6} - Dans l'exer. 3 donner le système des $h$, $x$, $y$ de Weyl, qui fournit des formules sans dénominateurs. Donner la démonstration Van der Waerden, (cf. Jacobson, bas p. 123 et p. 124).

\S\,4 - \uuline{Théorèmes d'existence}, (nouveau titre).

Dans ce \S\ on fait deux constructions distinctes, qui amènent presque aux mêmes résultats. On en laissera tomber une, d'où disparition des th. 3 et 4. Ceci amène au plan suivant :

1) Construction de $\underline{a}=\underline{x}+\underline{h}+\underline{y}$ ;

2) Définition des gros modules $E(\lambda)$ ; liberté de $\underline{x}$ et $\underline{y}$ ;

3) Construction de $\underline{g}$ comme quotient de $\underline{a}$ par les $\unsure{\Theta|\theta}_{ij}$ ; alors $\underline{g}=\underline{h}+\sum_{\lambda}\underline{g}^{\lambda}$, est de dim. finie, et est semi-simple ;

4) Définition des $W(\lambda)$ ; sont des $g$-modules ;

5) $W(\lambda)$ est de dim. finie si $\lambda$ est un pouah dominant ; conséquences pour $\underline{g}$.

Dans 3) la semi-simplicité de $\underline{g}$ se démontre ainsi : on a $[\underline{g}^{\alpha},\underline{g}^{-\alpha}]=\underline{h}^{\alpha}$, (vrai pour $\alpha$ simple, et $W$ opère) ; alors tout idéal commutatif $\underline{b}$ est nul, (car $\underline{b}\cap\underline{s}_\alpha=0$ pour toute sous-algèbre $\underline{s}_\alpha$ de dim. 3). A partir de là tous les cor. concernant $\underline{g}$ se démontrent facilement.
% archive sheet number: 191
\scanpage{193}
les $\xi_{\ell}$ du (i) dépendent du choix d'un pt \unsure{géom|géom.} de $X$, mais à isom non \unsure{uniques|unique} près (\unsure{définir|défin.|dépendance} \unsure{ens des|ens de|du} \uuline{$\unsure{t_{\alpha}|{\pi}_1}(X)$}] \inserted{[$\pi_1(X)$ opère sur $\xi_\ell$]} sont \unsure{indépendants} de ce choix.

Les $\xi_p^\alpha$ dépendent du choix d'un \deleted{corps} pt \inserted{de $X$} à valeurs dans un corps parfait sans plus, si on veut les $k(x)^{\unsure{\mathrm{prf}|\mathrm{pf}}}$\deleted{\unsure{,}} (clôture parfaite de $k(x)$), et pour $k$ donné, il ne \uline{semble pas} qu'il y ait \deleted{\unsure{aie}} des isom entre les $\xi_{p,a}^{(\alpha)}$ correspondant aux divers pts \unsure{à} valeurs dans $k$ (c'est la ``partie \unsure{continue}'' des $\unsure{H^1|H}$ p.ex. \unsure{qui va|que va} faire des ennuis), s'ils sont localisés en le $\unsure{=|\text{même}}$ $x\in X$, \deleted{\unsure{il est}} \deleted{\unsure{sont}} si $k$ est alg. clos, ils sont conjugués par autom. de $k$ \inserted{(étendus à $\unsure{\mathbb{W}|W}_p(k)$)} les groupes alg. associés sont donc conjugués sous autom de $k$ ---- il est $=$ gentil qu'ils veuillent être l'un tordu de l'autre \ldots\ Le groupe $\pi_1(X)$ n'opère plus sur $\xi_p^\alpha$,
% margin note written sideways
\margin{En revanche, on a comme structure supplémentaire les opérations $F,V$ avec $FV=p$}
par contre le groupe des autom de $k$ sur $k(x)$ \deleted{\unsure{[ou = de $k$|[ou = de $k$,}}
\deleted{\unsure{il existe]}} opère de façon $1/2$ \uline{linéaire} sur $\xi_p^{(\alpha)}$ ---- [\unsure{ou|on} \unsure{une|on} \unsure{donnée de|devient|devine}
% margin note written sideways
\margin{élégante. D'autre part, un autom $\sigma$ de $k$ qui n'est pas $k(x)$-linéaire transforme en général $\xi_{x,p}^{\unsure{\ast|\alpha}}$ en un $\xi_{xp}^{\sigma}$ qui n'est plus isom à $\xi_{x,p}$ ----}

Notations
\diagram{\(\begin{array}{ll}
\text{(i)}&
\left\{\begin{array}{ll}
\xi_{\bar{x},\ell} & \text{foncteur fibre $l$-adique en}\\
 & \text{le pt géom $\bar{x}$}\\
\xi_{x,\unsure{\ell|\beta}} & \text{par abus de langage pour $x\in X$}\\
 & \text{ou un pt de $X$ à val.}\\
 & \text{dans un corps quelc.}\\
 & \text{($\xi_{x,\ell}$ déf. à isom \uline{non unique} près)}
\end{array}\right.\\[1ex]
\text{(ii)}&
\left\{\begin{array}{ll}
\xi_{x,p}&\text{$x$ \deleted{\unsure{est}} pt de $X$ à val. dans un corps}\\
\xi_{x,p}^{0}&
\end{array}\right.
\end{array}\)}
% slanted
\margin{\uline{On peut mettre} \uline{des exposants $\alpha$} \uline{à volonté}}
% archive sheet number: 192
\scanpage{194}
ALGÈBRES de LIE SEMI-SIMPLES \mspage{(8)}\hfill Tribu n\textsuperscript{o} 59

Dire `` ces éléments $N_{\alpha\beta}$ '' à la p. 46, \unsure{l|1}.6 bas.

\uline{n\textsuperscript{o} 6} - Enoncer ainsi la prop. 2 : a) tout sous-espace $V$ stable par $\underline{h}$ est de la forme, (sous-espace $\underline{h}'$ de $\underline{h}$) $\oplus\sum_{\alpha\in\Sigma}\underline{g}^{\alpha}$ ; b) pour que $V$ soit une sous-algèbre, il faut et il suffit que $\Sigma$ soit clos et que $\alpha\in\Sigma\cap(-\Sigma)$ implique $H_\alpha\in\underline{h}'$. Exhiber alors le radical de $V$, (une partie vient de $\underline{h}$, une autre des $\alpha\in\Sigma$ tels que $-\alpha\notin\Sigma$).

Enoncé de la prop. 3 : $V$ réductive $\Longleftrightarrow\Sigma$ clos et symétrique ; $V$ semi-simple $\Longleftrightarrow\Sigma$\unsure{,|.} clos et symétrique et $\underline{h}'$ engendré par les $H_\alpha$ pour $\alpha\in\Sigma$ ; son système de racines est alors celui que dit la orop. 3.

Le cor. 2, p. \unsure{47|49}, ira en exer6. La dém. de la prop. 4 se simplifie, maintenant qu'on connait le radical. Dans le cor., p. 51, dire `` racines positives par rapport à une base de $R$ ''.

A propos des Borels, (p. 51 bas), dire que la forme quadratique est nulle sur $\underline{n}^{+}$ et $\underline{n}^{-}$, et que leurs orthogonaux sont $\underline{b}^{-}$ et $\underline{b}^{+}$. On parlera des paraboliques, (ou paraboréliques).

\uline{n\textsuperscript{o} 7} - Dans la prop. 7, indexer $S$, (qui s'appelle maintenant `` base de $R$ ''). A la l.8 bas, (p. 52), dire que les autres formules ``ont déjà été démontrées'', et donner les références.

Donner ici les relations $(\operatorname{ad}X_i)^{-a_{ij}+1}X_j=(\operatorname{ad}Y_i)^{-a_{ij}+1}Y_j=0$, (cf. \S\,4).

Pas de nom pour $(H_i,X_i,Y_i)$ en haut de la p. 53, ça conflicte ; on dira que ce système est associé à l'épinglage.

\uline{n\textsuperscript{o} 8} - Au début préciser que $X_\alpha\in\underline{g}^{\alpha}$. Le lemme 4 remonte au \S\,2, n\textsuperscript{o} 2. Introduire $p$ aussi dans l'énoncé de la prop. 6. La formule du bas de la p. 54 tient à ce que $\alpha+\beta$ est une racine ; le dire.
% archive sheet number: 193
